\section{Arithmetic and categorical reductions}\label{sec:reductions}
\subsection{Conventions and imported criteria}\label{sec:inputs}
Write $N_{ij}^{k}$ for the coefficient of $b_k$ in $b_i b_j$ and $i^*$
for duality. The matrices used in the software have entries
$(N_i)_{jk}=N_{ij}^{k}$. Consequently they give the right-handed matrix
convention
\[
N_iN_j=\sum_k N_{ji}^{k}N_k.
\]
This distinction is harmless for commuting matrices, but essential when
checking noncommutative signed tables. Frobenius reciprocity gives
$N_{i^*}=N_i^{\mathsf T}$.
The positive dimension character is denoted $d_i=\FPdim(b_i)$ and
$D=\sum_i d_i^2$.

The complexified based ring is a semisimple $*$-algebra. For an
irreducible representation $E$ of dimension $r_E$, let $f_E$ be its
formal codegree, in the convention of \cite{Ocode}. The regular
Casimir matrix
\[
C_R=\sum_i N_iN_{i^*}
\]
has eigenvalue $r_Ef_E$ with multiplicity $r_E^2$. It is only in the
commutative case that all these multiplicities are one and its
characteristic roots are themselves the formal codegrees.

We use the following established necessary conditions. Their precise
applications are certified by polynomial, prime, dimension-field, and
signed-table witnesses.

\begin{proposition}[Established inputs]\label{prop:inputs}
For a complex fusion categorification, the following hold.
\begin{enumerate}[label=(\roman*)]
\item Every formal codegree is a $d$-number \cite{Ocode}. If
$p(t)=t^s+a_1t^{s-1}+\cdots+a_s$ is its minimal polynomial, this is
expressed by $a_s^i\mid a_i^s$ for $1\le i\le s$.
\item For a commutative Grothendieck ring, character values are cyclotomic
integers. The polynomial factorization and finite-field witnesses used
here are those of \cite[Section~2]{DP4}; the commutativity hypothesis is
checked before this test is used.
\item A fusion subring gives a fusion subcategory. Its Frobenius--Perron
dimension divides that of the ambient category in the ring of algebraic
integers. A subring of rank two with $x^2=1+mx$, $m\ge2$, is excluded
by \cite{Ostrik2}.
\item Weakly integral fusion categories are pseudo-unitary
\cite{ENO}. For the two-Frobenius--Perron-dimension reduction we use the
Galois-conjugacy result of \cite[Lemma~3.1]{SchTwo}.
It is a reduction to a pseudo-unitary Galois conjugate, not a claim that
every originally given spherical structure is positive.
\item The zero-spectrum obstruction of \cite[Theorem~2.14]{LPR}
is a general categorification obstruction. Its finite hypothesis is
tested directly on the full ordered tensor.
\item The spherical dimension, normalized norm, sign-object moment,
adjoint-parity, and low-induction criteria are those of
\cite[Sections~2,3,6,7]{DP4}. The finite joint relative-trace criterion is the named theorem of
\cite{DP5}.
\end{enumerate}
\end{proposition}

For completeness, the finite tensor hypothesis in input (v) is as
follows. Put $\langle ab,cd\rangle=\sum_kN_{ab}^kN_{cd}^k$.
There must be nine indices $a,b,c,u,v,w,x,y,z$ such that
\[
 N_{ua}^{w}N_{vu}^{b}N_{vw}^{c}N_{xz}^{a}N_{bx}^{y}N_{yz}^{c}>0,
 \quad N_{ba}^{c}=1,
 \quad \sum_kN_{ux}^{k}N_{v^*y}^{k}N_{wz^*}^{k}=0,
\]
and each of the two sets
\[
 \{\langle vu,ca^*\rangle,\langle bu^*,cw^*\rangle,
       \langle v^*b,wa^*\rangle\},
 \qquad
 \{\langle bx,cz^*\rangle,\langle yx^*,ca^*\rangle,
       \langle b^*y,az^*\rangle\}
\]
contains $1$. Theorem~2.14 of \cite{LPR} excludes any such tensor.
The certificate stores the nine indices and the checker reevaluates all
these expressions. No associator is assumed to be unitary or pivotal.

Input (v) is especially inexpensive because it needs just a small
integer witness. We do not import the negative part of any integral
classification table or quadratic-family classification as a list of
forbidden labels.

\begin{corollary}[Certified lower-rank descent]\label{cor:subring}
Suppose a rank-six fusion ring contains a based fusion subring of rank
four and multiplicity at most $32$, or of rank five and multiplicity at
most $17$. If that subring is not based-isomorphic to one of the positive
rings in the respective main theorem of \cite{DP4,DP5}, the rank-six
ring is not categorifiable over $\CC$.
\end{corollary}
\begin{proof}
A hypothetical categorification would contain the fusion subcategory
on the prescribed basis subset. Apply the cited complete lower-rank
classification. The test enumerates unit-containing dual-closed subsets,
checks product closure, and compares the restricted tensor with every
positive model under every unit-fixing permutation. No negative labels
are read.
\end{proof}


\subsection{A certificate of a nonabelian character field}
\begin{theorem}[Noncommuting field automorphisms]\label{thm:automorphisms}
Let $R$ be a commutative fusion ring. Suppose an irreducible factor
$f\in\QQ[t]$ of the characteristic polynomial of some multiplication
matrix admits polynomials $a,b\in\QQ[t]$ such that
\[
 f(a(t))\equiv f(b(t))\equiv0\pmod f,
 \qquad a(b(t))-b(a(t))\not\equiv0\pmod f.
\]
Then $R$ has no complex fusion categorification.
\end{theorem}
\begin{proof}
In a categorification, the eigenvalues in question are values of
one-dimensional characters of $R\otimes\CC$ and are cyclotomic integers
by Proposition~\ref{prop:inputs}(ii). Hence $K=\QQ[t]/(f)$ embeds into an
abelian cyclotomic extension of $\QQ$. Every subfield of an abelian
Galois extension is Galois with abelian Galois group. Substitution by
$a$ and $b$ defines unital endomorphisms of the field $K$; they are
injective, hence surjective by finite dimensionality. The last displayed
condition says that these automorphisms do not commute, a contradiction.
\end{proof}

Factoring $f$ over $\QQ[t]/(f)$ is allowed in discovery, but is unnecessary
in verification: irreducibility, divisibility, two substitutions and one
nonzero remainder suffice. In the present census the test excludes
$\RR{2}{119}$, with the degree-six factor
\[
 f(t)=t^6-3t^5-2t^4+9t^3-5t+1.
\]
The two substitution polynomials and their nonzero commutator are stored
in the certificate and recomputed from this polynomial. This test is a
stronger certificate generator for the existing cyclotomicity condition,
not a new categorical hypothesis.

The same automated test excludes $\RR{6}{1891}$. A multiplication
matrix has the irreducible factor
\[
 t^6-11t^5+5t^4+38t^3-5t^2-11t-1,
\]
and the certificate supplies two noncommuting automorphisms of its
degree-six field. Their defining substitutions and nonzero commutator
are checked by exact polynomial reduction. Thus this case is decided
by the uniform arithmetic criterion without an induction search.

\subsection{Positive dimensions after a tensor-sign change}
\begin{lemma}[Tensor-sign Galois reduction]\label{lem:positive-sign}
Let $\mathcal C$ be spherical and let $\sigma$ be an embedding of its
simple-dimension field into $\mathbb R$. Put
$\eta_i=\operatorname{sgn}(\sigma(\dim X_i))$. If
\begin{equation}\label{eq:tensor-sign}
 \eta_0=1,\qquad N_{ij}^k>0\Longrightarrow\eta_i\eta_j=\eta_k,
\end{equation}
then $\mathcal C$ has a pseudo-unitary Galois conjugate after changing
its spherical structure by a tensor automorphism of the identity.
\end{lemma}
\begin{proof}
Extend $\sigma$ to the field of definition and conjugate the categorical
structure, as in the Galois reduction of \cite{DP4}. The scalars $\eta_i$
define a tensor automorphism of the identity precisely because of
\eqref{eq:tensor-sign}. They agree on duals and square to one, so twisting
the pivotal structure preserves sphericality. The resulting dimensions
are $|\sigma(\dim X_i)|>0$. A positive ring dimension character is the
Frobenius--Perron character, so this spherical structure is positive and
the conjugate is pseudo-unitary.
\end{proof}

This lemma must not be replaced by arbitrary basis sign changes: the
support condition \eqref{eq:tensor-sign} is essential. The spherical
checker now tests it at every real embedding of a remaining character
component. Combined with the complete signed-sphericalization tests of Section~\ref{sec:signed}, it supplies general exclusion certificates
for $\RR{2}{37}$, $\RR{2}{57}$ and $\RR{3}{204}$. The new step concerns
spherical character components; the exhaustive signed analysis is still
required before a conclusion is promoted to general nonexistence.

\subsection{Signed sphericalization at rank six}\label{sec:signed}
Let $\mathcal C$ be an arbitrary fusion category. Its canonical
sphericalization $\widetilde{\mathcal C}$ has a sign object
$\varepsilon$ of dimension $-1$ and two lifts of every simple object,
interchanged by tensoring with $\varepsilon$; see
\cite{ENO,Ostrik3}. Choose a unit-preserving set of lifts. Passing to
\[
T=K_0(\widetilde{\mathcal C})/(\varepsilon+1)
\]
gives integral signed structure constants satisfying
\begin{equation}\label{eq:signedbounds}
 |T_{ij}^{k}|\le N_{ij}^{k},\qquad
 T_{ij}^{k}\equiv N_{ij}^{k}\pmod2.
\end{equation}
They retain associativity, the unit, duality, and Frobenius reciprocity.
All these statements concern a necessary quotient of an actual
sphericalization; conversely, an admissible signed table need not have a
categorification.

The sphericalization arguments in \cite{DP4,DP5} choose the lift of each
self-dual object with positive second indicator. We keep this normalization.
In particular, the square-channel congruences are modulo four rather
than merely modulo two whenever that normalization applies. Non-self-dual
lifts are chosen compatibly with duality. These restrictions are part of
the cited normalization theorem, not heuristic reductions of a search.

\subsubsection{An exhaustive oriented parameterization}
Set $t_{abc}=T_{ab}^{c^*}$. In the absence of commutativity, identify
only the orbits generated by
\begin{equation}\label{eq:orbits}
(a,b,c)\longmapsto(b,c,a),\qquad
(a,b,c)\longmapsto(c^*,b^*,a^*).
\end{equation}
These are the cyclic and dual-reflection consequences of reciprocity.
An arbitrary transposition is not allowed. Unit and duality coefficients
are fixed first. Give each remaining orbit its interval in
\eqref{eq:signedbounds}, impose the applicable square-channel
congruences, and substitute into all ordered associativity equations.

\begin{proposition}[Exhaustive signed search]\label{prop:oriented}
The finite search just described contains the signed quotient of every
hypothetical categorification, after the indicated lift normalization.
If all its associative tables are rejected by necessary signed
sphericalization criteria, the original ring has no complex
categorification.
\end{proposition}
\begin{proof}
The quotient construction gives \eqref{eq:signedbounds} and
\eqref{eq:orbits}. Thus every quotient occurs in the finite product of
integer domains. Every pruning operation is a necessary consequence of
an equation over those domains. At a leaf the full ordered associativity
identities are checked. Exhausting the product therefore exhausts every
possible normalized quotient. A genuine category would supply one of
the remaining tables and one of its spherical dimension characters, so
rejection of all tables and all such characters is contradictory.
\end{proof}

The implementation checks machine-integer bounds before native
computation and does not interpret overflow or timeout as exhaustion.
The exact verifier reconstructs the parameter domains and the complete
list of surviving signed tables; it does not merely recheck a collection
of stored rejected examples.

\subsubsection{When the signed algebra must be commutative}
Let $d=(1,d_1,\ldots,d_5)$ be the dimension character of the selected
lifts, $K=\QQ(d_1,\ldots,d_5)$, and
$\delta=\sum_i d_i d_{i^*}$. Inflating a representation of $R$ to
$K_0(\widetilde{\mathcal C})$ doubles its formal codegree. Since the
sphericalization has dimension $2\delta$, central induction gives
\begin{equation}\label{eq:fieldconstraint}
 \delta/f_E\in\mathcal O_K,\qquad f_E\in K
 \quad(E\in\Irr(R\otimes\CC)).
\end{equation}
Indeed, \cite[Theorem~2.13]{Ostrik3} gives a central simple of dimension
$\delta/f_E$, and its restriction is an integral combination of the
selected lifts and their negatives. Its dimension is a nonzero algebraic
integer in $K$, proving both assertions.

\begin{proposition}[Rank-six commutative descent]\label{prop:descent}
Suppose $R$ is commutative of rank six and has a hypothetical
categorification. Its signed quotient $T\otimes\CC$ is commutative if
any of the following tensor-checkable conditions holds:
\begin{enumerate}[label=(\alph*)]
\item all six basis elements are self-dual;
\item exactly two basis elements are self-dual;
\item a formal codegree has degree greater than two;
\item the formal codegrees generate two different quadratic subfields;
\item zero has multiplicity less than four in
$\det(tI-C_R)\bmod2$.
\end{enumerate}
\end{proposition}
\begin{proof}
The coefficient-of-unit pairing of a signed quotient is still positive
definite on the distinguished basis after complex conjugation. Its
regular matrices are adjoint under duality. Hence the algebra is
semisimple. It has a one-dimensional dimension representation. If it
is noncommutative of dimension six, its only possible Wedderburn type is
\begin{equation}\label{eq:wedderburn}
 \CC\oplus\CC\oplus M_2(\CC).
\end{equation}
The Galois orbit of the dimension character has at most two elements,
so $[K:\QQ]\le2$. Equation~\eqref{eq:fieldconstraint} excludes
\eqref{eq:wedderburn} in cases (c) and (d).

In case (a), duality is the identity on the signed basis. Being also an
anti-automorphism, it forces commutativity. For (b), consider duality as
a complex-linear involutive anti-automorphism of
\eqref{eq:wedderburn}. The spherical dimension character is fixed, so
both one-dimensional blocks are fixed. On $M_2(\CC)$, an involutive
anti-automorphism is of orthogonal or symplectic type, with linear trace
$2$ or $-2$. The total trace is therefore $4$ or $0$. In the distinguished
basis it is the number of fixed basis elements, namely $2$, a
contradiction.

Finally the unique two-dimensional representation of
\eqref{eq:wedderburn} is Galois invariant. Its formal codegree is a
rational algebraic integer, hence an integer. The regular Casimir of
$T$ has the even eigenvalue $2f_E$ with multiplicity four. Thus its
characteristic polynomial is divisible by $t^4$ modulo two. But
\eqref{eq:signedbounds} implies $C_T\equiv C_R\pmod2$, contradicting
(e).
\end{proof}

This proposition is a reduction, not a noncategorification criterion by
itself. Once its hypotheses are verified, the commuting-character
methods of \cite{DP4,DP5} apply to every signed branch. Otherwise we use
Proposition~\ref{prop:oriented}. If a noncommutative table of type
\eqref{eq:wedderburn} remains, its one-dimensional characters are
computed as the common kernel of the matrix commutators. That space has
dimension two. The two-dimensional block and its formal codegree are
then recovered from the regular Casimir with the multiplicities in
\eqref{eq:wedderburn}. Each decomposition identity is checked exactly.

\subsection{A scoped literature obstruction at the first frontier}
\label{sec:premodular-frontier-bound}

The general existence question and the premodular existence question must
be kept separate. The following consequence of the complete rank-six
premodular fusion-rule catalogue is useful for the latter question.

\begin{proposition}\label{prop:rank6-premodular-dimension}
If a rank-six fusion ring admits a braided spherical categorification,
then its Frobenius--Perron dimension is at most $168$.
\end{proposition}
\begin{proof}
Use the fusion-rule classification of Dong--Schopieray
\cite[Section~5.3, Figures~13--15, 17 and 19]{DSNear}.
The integral representation-ring entries and their listed braided
variants have dimension at most $168$. The other entries belong to the
families below. Write $[j]_\ell=\sin(j\pi/\ell)/\sin(\pi/\ell)$.
The inequalities
\[
 \sin(\pi/\ell)\geq 2/\ell,\qquad
 |\sin(j\pi/\ell)|\leq1,\qquad
 \sqrt5<3,\qquad\sqrt{21}<5
\]
give the following rational upper bounds. Each is at most $168$.
\begin{center}
\begin{tabular}{lr}
\toprule
Family of fusion rules & Upper bound \\
\midrule
Pointed rank three times Fibonacci & $12$ \\
$\operatorname{Rep}(S_3)$ times Fibonacci & $24$ \\
Dimension-eight Ising and generalized-Ising rules & $8$ \\
Ising times Fibonacci & $16$ \\
Pointed rank two times $\mathcal C(A_1,7)_{\rm ad}$ & $51$ \\
Fibonacci times $\mathcal C(A_1,7)_{\rm ad}$ & $102$ \\
$\mathcal C(G_2,21)$ & $110$ \\
Dimension-twenty metaplectic rules & $20$ \\
$\mathcal C(A_1,13)_{\rm ad}$ & $2197/16$ \\
$\mathcal C(B_2,9)_{\rm ad}$ & $81$ \\
$\mathcal C(A_1,12)_{\rm ad}$ & $108$ \\
\bottomrule
\end{tabular}
\end{center}
For example, the Fibonacci dimension is less than $4$, and the sum of
squared dimensions of $\mathcal C(A_1,7)_{\rm ad}$ is at most
$1+2(7/2)^2$. For the two adjoint $A_1$ entries in the last three rows,
use their dimension formulae
$\tfrac{13}{4}\csc^2(\pi/13)$ and $3\csc^2(\pi/12)$.
For $B_2$, the formula is $9u^2$ with
$u=1-2\cos(8\pi/9)$ and $|u|\leq3$.
For $G_2$, its six dimensions are bounded by $1,4,4,4,5,6$.
These cover all rows of the cited catalogue. The caveat there about
counting braided equivalence classes for two families does not add
fusion rules to the list.
\end{proof}

The implementation derives the positive character and its number-field
embedding from each tensor. It then decides
$\operatorname{FPdim}(R)-168>0$ by exact real-root isolation, recording
the defining polynomial, dimension coordinates and isolating interval.
No tensor identifier is used to select the outcome. In particular,
the former multiplicity-three residual $\RR{3}{215}$ has the automatically
verified positive dimension vector
\[
 (1,1,d,d,d-1,d+1),\qquad d=4+\sqrt{17},
 \qquad\operatorname{FPdim}(R)=136+32\sqrt{17}>168.
\]
It therefore has no premodular categorification. Its general, pivotal,
spherical, pseudo-unitary and unitary existence questions are not
resolved by this proposition, because none of those hypotheses alone
supplies a braiding. They are resolved negatively by the independent
stabilizer-rotation argument of Theorem~\ref{thm:frontier-noncategorifiable}.

For a separate complete-tensor model audit, the rank-six ancillary data
of Ng--Rowell--Wen \cite{NRW12} contain $192$ modular records in $20$
Galois orbits and $11$ based-ring types. The supplied parser reads only
literal integer tensors and compares all $120$ permutations fixing the
unit. Their largest fusion coefficient is $2$, so this source supplies
no new model at multiplicity eight or for the multiplicity-three
frontier. This comparison is a model search within a specified
literature class, not an exclusion of arbitrary fusion categories.

\subsection{Biquadratic spherical traces}
\label{subsec:biquadratic-spherical-traces}
The following test uses sphericalization to obstruct a fusion category without
assuming that it itself admits a spherical structure.  Throughout this
subsection the input ring is commutative, self-dual, and of rank six.
Write $N_i$ for its regular multiplication matrices and
$T_0=\operatorname{Tr}(\sum_i N_i^2)$.

\begin{lemma}[Dimension-field containment]\label{lem:dimension-field-containment}
If the input ring is categorifiable, there is a commutative integral signed
based algebra $A$ of rank six, with symmetric multiplication matrices $M_i$,
such that $M_i\equiv N_i\pmod2$ and $|(M_i)_{jk}|\le (N_i)_{jk}$.
It has a dimension character $d=(1,d_1,\ldots,d_5)$ in a totally real field
$K$ of degree at most six.  If $\Delta=\sum_i d_i^2$, every original formal
codegree $f$ belongs to $K$, and $\Delta/f$ is an algebraic integer in $K$.
The same statement holds for every formal codegree of $A$.
\end{lemma}
\begin{proof}
Sphericalization has two lifts of each simple and a central sign object
$\epsilon$ of dimension $-1$.  Each lift of a self-dual simple is self-dual:
its two possible lifts have opposite nonzero dimensions, whereas a dual has
the same dimension.  Thus its Grothendieck ring is commutative.  Quotienting
by $\epsilon+1$ gives $A$ and the asserted bounds and parity.  Its regular
matrices are real symmetric, hence simultaneously diagonalizable with real
characters.  Evaluation at the dimension character is a field quotient
of $A\otimes\mathbb Q$, proving $[K:\mathbb Q]\le6$.
An original character extends to sphericalization with $\epsilon=1$ and
formal codegree $2f$.  The induction theorem for formal codegrees
\cite[Theorem~2.13]{Ostrik3} produces a central simple object of dimension
$2\Delta/(2f)=\Delta/f$.  Forgetting it to sphericalization expresses this
dimension as an integer combination of the numbers $\pm d_i$.
The same argument applies to signed characters with $\epsilon=-1$.
\end{proof}

Suppose the original formal-codegree fields generate
$E=\mathbb Q(\sqrt p,\sqrt q)$, where $p,q>1$ are squarefree, $p\equiv1\pmod4$,
and their quadratic field discriminants are coprime.  The lemma forces
$K=E$, since $4\mid[K:\mathbb Q]\le6$.  Write
$d_i=(\alpha_i+\beta_i\sqrt p+\gamma_i\sqrt q+
\delta_i\sqrt{pq})/4$.  Orthogonality of the four distinct Galois-conjugate
characters gives
\[
16+\sum_{i=1}^5\alpha_i^2
=p\sum_{i=1}^5\beta_i^2
=q\sum_{i=1}^5\gamma_i^2
=pq\sum_{i=1}^5\delta_i^2.
\]
Consequently $\operatorname{Tr}_{E/\mathbb Q}(\Delta)\in pq\mathbb Z_{>0}$.
The two remaining signed formal codegrees are greater than one, since all
are positive and their reciprocals sum to one.  Since the signed Casimir
trace is at most $T_0$,
\begin{equation}\label{eq:biquadratic-trace-bound}
\operatorname{Tr}_{E/\mathbb Q}(\Delta)<T_0-2.
\end{equation}

\begin{proposition}[Finite biquadratic obstruction]
\label{prop:biquadratic-trace-obstruction}
Assume in addition that the Frobenius--Perron formal codegree has minimal
polynomial $X^2-5nX+25n$, where $n=s^2+4=k^2\ell$ and
$\ell\in\{q,pq\}$.  Every categorification supplies an integral element
\[
\theta=\frac{\Delta}{5k\sqrt\ell}
=\frac{x+y\sqrt p+z\sqrt q+w\sqrt{pq}}4
\]
subject to all the following finite tests:
\begin{enumerate}
\item $\Delta$ is totally positive, its trace satisfies
\eqref{eq:biquadratic-trace-bound}, and $pq$ divides that trace.
Put $m=z$ if $\ell=q$ and $m=w$ if $\ell=pq$.  Then $m>0$ and
$\operatorname{Tr}(\Delta)=5k\ell m$.
\item $\operatorname{Norm}(\theta)$ divides $\ell^2m^4$.
\item Put $1-\operatorname{Tr}(1/\Delta)=A/B$ in lowest positive terms.
For some positive integer $j$, the last two signed codegrees have sum $Aj$
and product $Bj$, both roots greater than one, and sum at most
$T_0-\operatorname{Tr}(\Delta)$.  If their quadratic polynomial is
irreducible, $Bj\mid(Aj)^2$ and its quadratic field embeds into $E$.
\item The product of this quadratic polynomial and
$\prod_{\sigma:E\hookrightarrow\mathbb R}(X-\sigma\Delta)$ is congruent
modulo two to the original Casimir characteristic polynomial.
\end{enumerate}
If the finite set is empty, no complex fusion category realizes the ring.
\end{proposition}
\begin{proof}
For $u=(s+\sqrt n)/2$ one has $u^2-su-1=0$ and
$D=5\sqrt n\,u$.  Thus $\theta=u\Delta/D$ is integral by
Lemma~\ref{lem:dimension-field-containment}.
The integral basis is the product of the integral bases of the two quadratic
fields, because their discriminants are coprime.  If $(e,h)=(1,2)$ for
$q\equiv1\pmod4$, and $(e,h)=(0,1)$ otherwise, its congruences are
\[
\begin{aligned}
hw&\equiv0\pmod2,& h(z-w)&\equiv0\pmod4,\\
y-ew&\equiv0\pmod2,& x-y-ez+ew&\equiv0\pmod4.
\end{aligned}
\]
Total positivity supplies the strict coordinate bounds
$|w|<m/\sqrt p$, $|x|<m\sqrt q$, $|y|<m\sqrt{q/p}$ when $\ell=q$;
and $|x|<m\sqrt{pq}$, $|y|<m\sqrt q$, $|z|<m\sqrt p$ when $\ell=pq$.
Together with the trace bound these give the finite search.
The global dimension is a $d$-number by \cite[Theorem~1.2]{Ocode}.  For a nonzero $d$-number $a$ in a
Galois field, $\operatorname{Tr}(a)^4/\operatorname{Norm}(a)$ is integral:
$\operatorname{Tr}(a)/a$ is integral and one takes its norm.  Applying this
to $\Delta$ gives the norm divisibility in (2).
The two complementary codegrees are Galois stable, so their sum and product
are integers.  The reciprocal identity gives $(Aj,Bj)$; the quadratic
$d$-number test and Lemma~\ref{lem:dimension-field-containment} give (3).
Finally $M_i\equiv N_i\pmod2$ implies the congruence in (4).
\end{proof}

The tensor-only implementation \texttt{new\_biquadratic\_trace.py} excludes
\texttt{R6\_7\_3813}, \texttt{R6\_8\_6940}, \texttt{R6\_8\_6973}, and
\texttt{R6\_8\_6998}.  Their fields are respectively
$\mathbb Q(\sqrt5,\sqrt{193})$,
$\mathbb Q(\sqrt5,\sqrt{34})$,
$\mathbb Q(\sqrt{17},\sqrt{257})$, and
$\mathbb Q(\sqrt5,\sqrt{13})$.
For \texttt{R6\_8\_6973}, the trace is $2570m$, the trace bound gives
$1\le m\le2$, and the lattice condition requires $17\mid m$.
For \texttt{R6\_8\_6998}, the only remaining complementary pair has polynomial
$X^2-2350X+2500$; its product with the four dimension codegrees is
$X^6$ modulo two, whereas the original Casimir polynomial is $X^6+X^2$.
All four certificates replay with Python optimization enabled.  The checker
also reproduces the imported multiplicity-thirteen example's 1129 totally
positive candidates, rejects altered certificates, respects simultaneous
basis permutations, and excludes none of the 54 established positive models.

\subsection{Formal-codegree trace ideals}
\label{subsec:formal-codegree-trace-ideals}

The field-containment lemma admits a useful consequence even when the signed
dimension field has not been identified.  For an algebraic integer $f$, write
\[
\mathfrak t(f)=\operatorname{Tr}_{\mathbb Q(f)/\mathbb Q}
       \bigl(f\mathcal O_{\mathbb Q(f)}\bigr)=\tau(f)\mathbb Z,
\qquad \tau(f)>0.
\]
Conjugating $f$ does not change this integer ideal.

\begin{theorem}[Formal-codegree trace-ideal obstruction]
\label{thm:formal-codegree-trace-ideal}
Let $R$ be a commutative self-dual rank-six fusion ring, with regular matrices
$N_i$ and Casimir trace $T_0=\operatorname{Tr}(\sum_iN_i^2)$.
Choose any collection of its formal codegrees $f$ and put
$L=\operatorname{lcm}_f\tau(f)$.
If $R$ is categorifiable, the signed sphericalization has a dimension field
$K$ and global dimension $\Delta$ satisfying
\[
0<\operatorname{Tr}_{K/\mathbb Q}(\Delta)\in L\mathbb Z,
\qquad
\operatorname{Tr}_{K/\mathbb Q}(\Delta)\le T_0.
\]
In particular, $L>T_0$ excludes every complex fusion categorification.
No pivotality assumption or local-field isolation is required.
\end{theorem}
\begin{proof}
Lemma~\ref{lem:dimension-field-containment} gives $F=\mathbb Q(f)\subseteq K$
and $\Delta/f\in\mathcal O_K$.  The relative trace of an algebraic integer
is integral.  Therefore
\[
\operatorname{Tr}_{K/F}(\Delta)
=f\operatorname{Tr}_{K/F}(\Delta/f)\in f\mathcal O_F.
\]
Transitivity of trace gives the asserted divisibility for every selected $f$.
The embeddings of $K$ give distinct characters of the signed algebra, with
formal codegrees the corresponding conjugates of $\Delta$.  All its formal
codegrees are positive.  Thus their total, the signed Casimir trace, is at
least $\operatorname{Tr}_{K/\mathbb Q}(\Delta)$.  The signed coefficient
bounds give
$\operatorname{Tr}(\sum_iM_i^2)=\sum_{i,j,k}(t_{ij}^k)^2
\le\sum_{i,j,k}(N_{ij}^k)^2=T_0$.
\end{proof}

The implementation uses every rational and quadratic original codegree.
For a rational integer $f$, $\tau(f)=f$.  For a quadratic codegree
$f=(A+b\sqrt q)/2$, with $q>1$ squarefree, the integral basis gives
\[
\tau(f)=
\begin{cases}
\gcd\bigl(A,(A+bq)/2\bigr),&q\equiv1\pmod4,\\
\gcd(A,bq),&q\equiv2,3\pmod4.
\end{cases}
\]
Omitting higher-degree factors weakens the test and creates no exclusion.
The calculation depends only on the complete multiplication tensor.

For \texttt{R6\_8\_6997}, the rational formal codegrees include $15$, while
the FP codegree has minimal polynomial $X^2-6145X+30725$ and trace ideal
$6145\mathbb Z$.  Hence $L=18435>T_0=6173$, giving a general exclusion.
The same test supplies independent exclusions of
\texttt{R6\_7\_3812}, \texttt{R6\_8\_6234}, \texttt{R6\_8\_6509},
\texttt{R6\_8\_6940}, and \texttt{R6\_8\_6973}.

\begin{corollary}[Signed saturation from a trace ideal]
\label{cor:trace-ideal-saturation}
For an unordered triple of nonunit basis labels with coefficient $a\ge2$,
let $w$ be its number of ordered permutations, and put
\[
\eta=\min w\bigl(a^2-(a-2)^2\bigr),
\]
where the minimum runs over those triples.  If $0\le T_0-L<\eta$, every
signed sphericalization coefficient has absolute value equal to its
original nonnegative coefficient.  All remaining choices are finite sign
choices subject to associativity.
\end{corollary}
\begin{proof}
Parity forces any decrease of absolute value to be at least two.  Full
symmetry forces the same decrease on all $w$ ordered permutations.  Such a
decrease loses at least $\eta$ from the Casimir trace, whereas
Theorem~\ref{thm:formal-codegree-trace-ideal} permits loss at most $T_0-L$.
The unit coefficients do not change in the signed quotient.
\end{proof}

For \texttt{R6\_8\_6520}, one has $L=6145$, $T_0=6165$, and $\eta=24$.
Thus its signed quotient is saturated.  Exhaustive exact signed associativity
leaves one table up to the certified sign gauges; that table is gauge-equivalent
to the original ring.  The non-FP character orbits have zero entries and cannot
be dimension characters.  Consequently every categorification would have a
pseudo-unitary Galois conjugate.  Proposition~\ref{prop:quadratic-gauss-family}
excludes that possibility, completing the general noncategorifiability proof.

The authoritative implementation is \texttt{new\_trace\_ideal.py}.  It checks
all arithmetic from the tensor, replays signed enumeration when used, and
rejects altered certificates and omitted signed tables.  Its control report
uses the frozen initial forty residual tensors and all fifty-four established
positive models; the latter are all admitted by this necessary test.

\subsection{An all-orders quadratic Gauss obstruction}
\label{subsec:quadratic-gauss-family}
The trace-ideal reduction also gives a complete exclusion at
multiplicity eight.  The argument below treats every possible central
restriction column and every root-of-unity order; it does not impose a
numerical conductor cutoff.

Write $R(m)$ for the commutative self-dual rank-six tensor on
$1,a,b,c,d,e$ whose nonunit coefficients equal $m$, except that the four
unordered triples
\[
 (a,a,b),\ (a,a,c),\ (a,b,b),\ (b,b,c)
\]
have coefficient $m-1$, and
\[
 (c,d,d),\ (c,e,e),\ (d,d,e),\ (d,e,e)
\]
have coefficient $m+1$.  Coefficients are invariant under permutations of
the three indices, and the unit coefficients are fixed by self-duality.
Set
\[
 n=25m^2+4,\quad U=\frac{5m+\sqrt n}{2},\quad
 j=(1,1,1,1,1)^t,\quad\delta=(-1,-1,0,1,1)^t.
\]
The FP dimensions are $(1,U-1,U-1,U,U+1,U+1)$, and
$D=5(U^2+1)=(5n+25m\sqrt n)/2$.
The other five characters are the quadratic conjugate of the FP
character and, with $p=(1+\sqrt5)/2$, $q=(1-\sqrt5)/2$,
\[
 (1,p,q,-1,0,0),\ (1,q,p,-1,0,0),\
 (1,0,0,1,-p,-q),\ (1,0,0,1,-q,-p).
\]
The corresponding formal codegrees are $D,D',5,5,5,5$.
These identities, character completeness, and the full tensor are checked
exactly by \filename{software/new\_gauss\_family.py}.

\begin{proposition}[Quadratic Gauss family obstruction]
\label{prop:quadratic-gauss-family}
The tensor $R(7)=\RR{8}{6520}$ has no pseudo-unitary categorification.
Consequently the trace-ideal signed reduction excludes every complex
fusion categorification of this tensor.
More generally, the same exact verifier gives a necessary criterion for
$R(m)$, $m\ge3$: whenever its exhaustive column enumeration has the
families described below, it requires
\[
 12m^2+2\ \ge\ \frac{\varphi(f)m b}{2c},\qquad
 n=b^2q_0,\quad f=\operatorname{disc}\QQ(\sqrt{q_0}),\quad
 c=\sqrt{f/q_0},
\]
where $q_0$ is squarefree.
\end{proposition}
\begin{proof}
Assume a pseudo-unitary categorification and use its positive spherical
structure.  If $v_Y$ is the restriction column of a central simple, then
$\sum_Yv_Yv_Y^t=G$, where $G=\sum_iN_i^2$.
The unit column of $G$ is $(6,5mj+\delta)^t$, and its nonunit block is
\[
 5I+(25m^2+1)jj^t+5m(j\delta^t+\delta j^t)+\delta\delta^t.
\]
By central induction, $I(1)$ contains six distinct simples with dimensions
$1,z=U^2,t=U^2+1,t,t,t$ and twists one
\cite[Theorems 2.5 and 2.13]{Ostrik3}.
The symmetry of the central $S$-matrix gives, for character-labelled low
columns,
\[
 \dim(Y_\chi)\chi(v_\psi)=\dim(Y_\psi)\psi(v_\chi).
\]
Applying this with the quadratic-conjugate character and the four
codegree-five characters determines the $z$ column as $(1,mj)^t$:
indeed the conjugate dimension of a $t$ column is $t'=t/z$, so
reciprocity gives each small character value as $zt'/t=1$.  The complete
character table then determines the column uniquely.  Write the remaining four columns as
$(1,mj+w_\nu)^t$.  Their dimensions and the unit column of the Gram identity
give
\[
 j^tw_\nu=0,\qquad \delta^tw_\nu=1,\qquad
 \sum_\nu w_\nu=\delta.
\]
On the three-dimensional subspace perpendicular to $j$ and $\delta$,
$G$ restricts to $5I$.  Taking the trace of its residual Gram inequality
therefore gives $\sum_\nu\|w_\nu\|^2\le16$.  Each squared norm is at least
two, so every $w_\nu$ has squared norm at most ten and coordinates in
$[-3,3]$.

The exhaustive enumeration has 32 individual deviations, 59 unordered
quadruples satisfying the sum and norm conditions, and 11 satisfying a
labelling by the four small characters and modular reciprocity.
Exactly three have positive-semidefinite residual Gram matrix.  Their
rows of deviations are
\[
\begin{split}
W_1&=\begin{pmatrix}-1&0&1&-1&1\\-1&1&-1&1&0\\0&-1&1&1&-1\\1&-1&-1&0&1\end{pmatrix},\qquad
W_2=\begin{pmatrix}-1&0&1&0&0\\0&-1&1&0&0\\0&0&-1&0&1\\0&0&-1&1&0\end{pmatrix},\\
W_3&=\begin{pmatrix}-1&0&1&1&-1\\-1&1&-1&0&1\\0&-1&1&-1&1\\1&-1&-1&1&0\end{pmatrix}.
\end{split}
\]
All enumeration bounds have just been proved; no low column is selected
in advance.

A remaining column is $(0,v)^t$, with $v\in\ZZ_{\ge0}^5\setminus\{0\}$.
Let $Q$ be the nonunit residual Gram matrix.  The necessary conditions are
\[
 Q-vv^t\succeq0,\qquad Q-vv^t\ge0\text{ entrywise},\qquad
 D/\dim(v)\text{ algebraically integral}.
\]
The last condition follows because $D^2/\dim(v)^2$ is a formal codegree
of the modular center.  All three residuals annihilate
$(1,1,-4,1,1)^t$.  Hence the first, second, fourth, and fifth coordinates
determine the third.  The bounds $v_i\le\lfloor\sqrt{Q_{ii}}\rfloor$ and
\[
 |v_a-v_b|\le\lfloor\sqrt{Q_{aa}+Q_{bb}-2Q_{ab}}\rfloor,
 \qquad
 |v_d-v_e|\le\lfloor\sqrt{Q_{dd}+Q_{ee}-2Q_{de}}\rfloor
\]
make the enumeration finite and complete.  Gram positivity is checked
with exact rational Schur complements, or equivalently by the kernel
condition and $v_J^tQ_J^{-1}v_J\le1$ on a full-rank principal submatrix.
For a quadratic number $x+y\sqrt n$, integrality is equivalent to integral
trace $2x$ and integral norm $x^2-ny^2$.

For $m=7$ the counts at the bounded, Gram-compatible, and integral-quotient
stages are respectively
\[
 (511,58,1),\qquad (6147,282,9),\qquad (511,58,1).
\]
The outer branches admit only $j$, although $Q$ has rank two, and are
therefore impossible.  The complete middle list is
\[
 j,\quad j\pm(e_a-e_b),\quad j\pm(e_d-e_e),\quad
 7j-e_i+e_k\quad(i\in\{a,b\},\ k\in\{d,e\}).
\]
For the general verifier, the accepted list may also include the first
five types with $j$ replaced by $2j$ and the last four types with $m$
replaced by $2m$; any additional high type causes the criterion to abstain.

Put $u=e_a-e_b$, $v=e_d-e_e$.  Exact Gram identities are
\[
 \delta^tQ\delta=32,\quad\delta^tQj=80m,\quad
 u^tQu=v^tQv=8,\quad
 \delta^tQu=\delta^tQv=u^tQv=0.
\]
If $n_m,n_{2m}$ count the two possible cross-column families, the first
two equations give $n_m+n_{2m}=8$ and $n_m+2n_{2m}=8$.
Thus exactly eight cross columns occur, all with coefficient $m$.
They saturate both $u,v$ squared forms, excluding all offset types.
The three vanishing mixed forms force each of the four columns
$B_{ik}=mj-e_i+e_k$ to occur twice.  Subtracting them leaves
\[
 Q-2\sum_{i,k}B_{ik}B_{ik}^t=(12m^2+2)jj^t.
\]
Thus the number of uniform columns, with a column $2j$ counted four
times, is $L=12m^2+2$.

Let $\Theta_{ik}$ be the sum of the two twists belonging to $B_{ik}$,
and let $\Lambda$ be the sum of uniform twists with these weights.
It is a sum of exactly $L$ roots of unity.  The first twist moment of
$I(X)$ vanishes for nonunit $X$ \cite[Theorem 2.5]{Ostrik3}, so, putting
$w=5U$, the five equations combine as
\[
 w\Lambda j+t\sum_{i,k}\Theta_{ik}B_{ik}
 =-mzj-t(4mj+\delta).
\]
Pairing with $\delta$ gives $\sum\Theta_{ik}=-2$; pairing with $j$ gives
\[
 w\Lambda=-mz-2mt,\qquad
 \Lambda=-\frac{m^2+m\sqrt n}{2}.
\]

For completeness the required root-of-unity bound holds for all orders.
For $K=\QQ(\sqrt{q_0})$ with positive fundamental discriminant $f$, let
$c=\sqrt{f/q_0}$.  The coefficient of $\sqrt{q_0}$ in the normalized trace
to $K$ of a single root of unity has absolute value at most
$c/\varphi(f)$.  If the root's cyclotomic field does not contain $K$,
the trace is rational.  Otherwise its order is divisible by $f$.
Averaging over the kernel of the primitive quadratic character expresses
the irrational part as an induced quadratic Gauss sum divided by
$\varphi(N)$.  This sum is zero unless the extra factor $N/f$ is squarefree
and coprime to $f$; in the nonzero case its magnitude is $\sqrt f$.
The formula follows directly by Chinese remaindering: excess exponents
at primes of $f$ cancel over lifts, and the coprime factor contributes its
Ramanujan sum.  As $\varphi(N)\ge\varphi(f)$, the stated bound follows.
Normalized trace and the triangle inequality then require
\[
 L\ge\frac{\varphi(f)}c\left|[\sqrt{q_0}]\Lambda\right|
   =\frac{\varphi(f)mb}{2c}.
\]
For $m=7$, $n=q_0=f=1229$, $b=c=1$, and $\varphi(f)=1228$.
The forced sum of 590 roots is $-(49+7\sqrt{1229})/2$, but its required
length is at least $1228\cdot7/2=4298>590$.  This contradiction proves the
pseudo-unitary exclusion.  The separate complete signed trace-ideal
certificate supplies a pseudo-unitary Galois conjugate of any hypothetical
categorification, proving the general exclusion.
\end{proof}

The executable certificate is tensor-based and invariant under relabelling
of the five nonunit objects.  Its normal and optimized-Python replay
reconstructs every low and high column, rejects missing branches and
altered arithmetic.  The frozen
40-input continuation scan applies the criterion to every input; only
$\RR{7}{3582}$ and $\RR{8}{6520}$ match this tensor family, and both are
excluded in the pseudo-unitary scope. The independent trace-ideal
reduction makes the latter exclusion unrestricted. The scan
of all $16612$ tensors through multiplicity eight checks the parameter-three through parameter-seven
members uniformly:
\[
\begin{array}{c|r|r}
\text{tensor}&\text{forced root count}&\text{Gauss lower bound}\\\hline
\RR{4}{803}&110&342\\
\RR{5}{1500}&194&400\\
\RR{6}{2554}&302&1440\\
\RR{7}{3582}&434&1344\\
\RR{8}{6520}&590&4298
\end{array}
\]
The table records a separate application of the same necessary criterion
to all five rings.

\subsection{Ramification and rational-codegree trace bounds}
\label{subsec:ramification-trace}
This criterion applies to arbitrary complex fusion categorifications.  It
does not assume a pivotal structure and does not enumerate selected
dimension characters.  Let $R$ be commutative and self-dual of rank six,
and write $T_0=\sum_{i,j,k}(N_{ij}^k)^2$.

\begin{proposition}[Ramified dimension field and Casimir trace]
\label{prop:ramification-trace}
Suppose the FP global dimension is
$D=a+b\sqrt q$, where $q>1$ is squarefree,
$q\equiv2,3\pmod4$, and $a,b\in\ZZ$, $b>0$.
Suppose also that
\[
 R/2R\cong\FF_2[z]/(z^2)\times\prod_j\FF_{2^{e_j}}.
\]
Let $\mathcal L\subseteq\ZZ^2$ be the lattice of pairs $(S,B)$ for which
\[
 \frac{2qaS-2qbB}{a^2-qb^2},\qquad
 \frac{2aB-2qbS}{a^2-qb^2},\qquad
 \frac{2qS}{c},\qquad\frac{2B}{c}
 \quad\hbox{are integers}
\]
for every rational original formal codegree $c$.
Let $v$ be the positive generator of the projection of $\mathcal L$
onto its first coordinate.  If
\[
 4qv+5>T_0,
\]
then $R$ has no complex fusion categorification.
\end{proposition}
\begin{proof}
Sphericalize a hypothetical categorification and quotient its Grothendieck
ring by the central sign object plus the unit.  The result is a rank-six
integral signed algebra $A$ with multiplication coefficients $t_{ij}^k$
satisfying
\[
 t_{ij}^0=\delta_{ij},\qquad
 |t_{ij}^k|\le N_{ij}^k,\qquad
 t_{ij}^k\equiv N_{ij}^k\pmod2.
\]
The two lifts of each original self-dual simple have opposite nonzero
spherical dimensions, so duality cannot interchange them.  Consequently
the regular matrices $M_i$ of $A$ are real symmetric and commute.
The dimension character has values $d_i$ in a totally real field factor
$K$ of $A\otimes\QQ$, and $\Delta=\sum_i d_i^2$ is the global dimension
of the original category.

For every original formal codegree $f$, its pulled-back character on
sphericalization has formal codegree $2f$.
The central-induction theorem
\cite[Theorem~2.13]{Ostrik3} therefore gives a central simple of dimension
$\Delta/f$.  Forgetting it expresses this dimension as an integer
combination of the $\pm d_i$.  Thus $\Delta/f\in\mathcal O_K$, and in
particular $D\in K$ and $\QQ(\sqrt q)\subseteq K$.
This also recalls the dimension-field argument of
Lemma~\ref{lem:dimension-field-containment}.

Since $A/2A=R/2R$, lift the displayed primitive idempotents to
$A\otimes\ZZ_2$.  A finite free factor with finite-field reduction is
finite etale over $\ZZ_2$, so its generic fiber is unramified.  The only
remaining factor has rank two.  But $\QQ_2(\sqrt q)$ is ramified for
$q\equiv2,3\pmod4$, and every component of $K\otimes\QQ_2$ contains
this field.  No such component can be supported on an unramified factor.
Hence $[K:\QQ]\le2$, and therefore $K=\QQ(\sqrt q)$.

Now $\mathcal O_K=\ZZ[\sqrt q]$.  Write
$d_i=\alpha_i+\beta_i\sqrt q$ for $1\le i\le5$, and set
$S=\sum_i\beta_i^2>0$ and $B=\sum_i\alpha_i\beta_i$.
The two conjugate dimension-character vectors are orthogonal, because
the $M_i$ are commuting real symmetric matrices.  It follows that
\[
 1+\sum_i\alpha_i^2=qS,\qquad
 \Delta=2qS+2B\sqrt q,\qquad
 \Delta+\Delta'=4qS.
\]
Integrality of $\Delta/D$ gives the first two conditions defining
$\mathcal L$.  Integrality of $\Delta/c$ gives the other two for every
rational codegree $c$.  Therefore $S\ge v$.

The signed Casimir $\sum_iM_i^2$ has the six signed formal codegrees as
its eigenvalues.  The two dimension characters contribute $4qS$.
Each of the other four codegrees is strictly greater than one: equality
would force every nonunit character value to vanish, contrary to the
unit coefficient in $x_i^2$.  Since the trace is integral,
\[
 4qv+5\le\Tr\Bigl(\sum_iM_i^2\Bigr)
 =\sum_{i,j,k}(t_{ij}^k)^2\le T_0,
\]
contradicting the asserted strict inequality.
\end{proof}

\begin{corollary}
The rings $\RR{7}{3812}$, $\RR{8}{6234}$, and $\RR{8}{6509}$ have no
complex fusion categorifications.
\end{corollary}
\begin{proof}
For all three rings,
\[
 q=226,\qquad D=2260+150\sqrt{226},\qquad
 D^2-4520D+22600=0.
\]
The following table lists their exact local factors and trace data.
Here $B_2=\FF_2[z]/(z^2)$ and $D'$ is the conjugate of $D$.
\begin{center}
\begin{tabular}{@{}llllr@{}}
\toprule
Ring & Algebra modulo two & Formal codegrees & $v$ & $T_0$\\
\midrule
$\RR{7}{3812}$ & $B_2\times\FF_4\times\FF_2^2$
 & $D,D',3,15,5,5$ & 15 & 4548\\
$\RR{8}{6234}$ & $B_2\times\FF_4\times\FF_2^2$
 & $D,D',3,15,5,5$ & 15 & 4548\\
$\RR{8}{6509}$ & $B_2\times\FF_2^4$
 & $D,D',3,3,15,15$ & 15 & 4556\\
\bottomrule
\end{tabular}
\end{center}
In particular $\Delta/15\in\mathcal O_K$ and
$\Delta=452S+2B\sqrt{226}$ imply $15\mid S$, because
$\gcd(452,15)=1$.  Thus the signed Casimir trace is at least
$904\cdot15+5=13565$, exceeding each displayed $T_0$.
\end{proof}

The finite local checks exhaust all $64$ elements modulo two, find every
primitive idempotent, and verify orthogonality and completeness of the
factor decomposition.  Every nonzero element of a claimed field factor
is checked to have an inverse; the remaining rank-two factor has exactly
one nonzero square-zero element.  The lattice step is recovered from exact
congruences, not floating-point bounds.  The three certificates require
neither a signed-table search nor a pseudo-unitary exclusion.

The tensor-based implementation is
\filename{software/new_ramified_signed.py}.  It was applied to the frozen
second-round frontier of $40$ rings, including all $39$ with multiplicity
four through eight.  Of those $39$, the hypotheses fail for $34$, one
($\RR{7}{3582}$) proceeds to a bounded signed reduction, and four are
excluded directly: the three above and an independent field-containment
exclusion of $\RR{8}{6940}$.  The latter is also covered by the
biquadratic criterion, so it is not counted twice.
The complete per-input report is
\filename{research/round2_higher/ramification_application.json}.
All $54$ positive models are preserved.  Ordinary and optimized Python
checks verify the three certificates after nontrivial basis permutations
and reject six certificate corruptions and a corrupted tensor.


\subsection{Local character fields and quadratic codegree saturation}
\label{sec:r3-local-fields}

The next criterion strengthens the signed Casimir bounds by determining
the dimension field before enumerating possible formal codegrees.
Write $A$ for the signed algebra supplied by
Lemma~\ref{lem:dimension-field-containment}, $K$ for its dimension
character field, and $\Delta$ for its dimension codegree. Thus every
original or signed formal codegree $f$ lies in $K$, and
$\Delta/f$ is an algebraic integer. In particular this statement
concerns the sphericalization of an arbitrary categorification; it
does not assume a pivotal structure on that category.

\begin{proposition}[Two local ways to isolate the dimension field]
\label{prop:r3-local-isolation}
Let $R$ be commutative and self-dual of rank six, and suppose an
original formal codegree generates $E=\mathbb Q(\sqrt q)$, where
$q>1$ is squarefree. Either of the following conditions implies
$K=E$ and that the other four signed formal codegrees are rational
integers greater than one:
\begin{enumerate}
\item $q\equiv5\pmod8$, and $R/2R$ is a product of finite fields
with exactly one even-degree factor, of degree two;
\item $R/2R$ has three local factors: $\mathbb F_8$, a rank-two
local algebra, and $\mathbb F_2$.
\end{enumerate}
In the second case the conclusion uses cyclotomicity of characters
of a categorifiable fusion ring.
\end{proposition}
\begin{proof}
The dimension-field containment lemma gives $E\subset K$.
Because the regular matrices of $A$ are symmetric, $A\otimes\mathbb Q$
is a product of totally real number fields. Its primitive local
idempotents at two lift from $A/2A=R/2R$ to $A\otimes\mathbb Z_2$.
A finite free local factor with finite-field special fibre is finite
\'etale: its relative differentials vanish by Nakayama, and it is
flat. It is therefore the unramified extension with that residue
degree. Equivalently one may apply the classification of finite
\'etale algebras over a henselian local ring; see the Stacks Project,
\href{https://stacks.math.columbia.edu/tag/04GK}{Tag 04GK}.

In the first case $E\otimes\mathbb Q_2$ is the unramified quadratic
field. Every completion of a number field containing $E$ has even
degree over $\mathbb Q_2$. Since $K$ is a field quotient of
$A\otimes\mathbb Q$, its completions must use only the unique
even-degree local factor. That factor has degree two, so $K=E$.
If a codegree belonging to another character field $F$ were
irrational, it would generate $E$, because it already belongs to
$K=E$. Then $F$ would require another even-degree local factor,
which is impossible. All remaining codegrees are therefore rational.

In the second case the lifted rank-three factor has generic fibre
the unramified cubic extension of $\mathbb Q_2$. Let $F$ be the
global character field containing this local factor. A signed
character extends to a character of the sphericalization with
the sign object sent to $-1$. Hence its values are cyclotomic by
\cite[Theorem~8.51]{ENO}; in particular $F/\mathbb Q$ is abelian
and Galois. All its completions at two have the same local degree,
here three. Thus $[F:\mathbb Q]$ is three or six. Degree six would
require two degree-three local factors, whereas the remaining
lifted blocks have ranks two and one. Hence $F$ is cubic and uses
the entire rank-three block. The remaining rational algebra has
dimension three. The field $K$ contains the quadratic field $E$,
so it lies in this remainder and has degree two. Thus $K=E$.
The other character fields are $F$ and $\mathbb Q$; their
intersections with $E$ are rational. The remaining codegrees,
which lie in these intersections, are rational as claimed.

In both cases rational formal codegrees are algebraic integers.
They are positive, and the sum of reciprocals of the six signed
codegrees is one. Each is consequently greater than one.
\end{proof}

The second argument does not assert that the whole two-adic order
is \'etale. In particular a rank-two order with dual-number
reduction may have split generic fibre. Only its rank and its
separation from the cubic block are needed.

\begin{proposition}[Finite quadratic spectrum bound]
\label{prop:r3-quadratic-spectrum}
Suppose $K=E=\mathbb Q(\sqrt q)$ and the complementary four
codegrees are rational, as in Proposition~\ref{prop:r3-local-isolation}.
Assume $q\equiv1\pmod4$ and the original Frobenius--Perron
codegree has the form $D=c\sqrt q\,u$, where $c$ is a positive
integer and $u$ is an integral unit of norm $-1$.
Let $T=\operatorname{Tr}(\sum_iN_i^2)$.
Every categorification supplies integers $x,y$ and
$f_1\leq\cdots\leq f_4$ satisfying the following finite system:
\[
\begin{gathered}
 \Delta=\frac{c(qy+x\sqrt q)}2,\qquad
 x\equiv y\pmod2,\quad y>0,\quad x^2<qy^2,\quad cqy\leq T-5,\\
 n=qy^2-x^2,\qquad n\mid4qy^2,\qquad
 \sum_{i=1}^4\frac1{f_i}=1-\frac{4y}{cn}>0,\\
 f_i>1,\qquad
 f_i\mid\gcd\left(\frac{c(qy-x)}2,cx\right),\qquad
 cqy+\sum_{i=1}^4f_i\leq T.
\end{gathered}
\]
In addition $\Delta/f$ is integral for every original rational
codegree $f$, and the polynomial
\[
 \left(z^2-cqy\,z+\frac{c^2qn}{4}\right)
       \prod_{i=1}^4(z-f_i)
\]
is congruent modulo two to the original Casimir polynomial.
If the system is empty, $R$ is not categorifiable. Otherwise its
least total trace bounds the possible signed Casimir deficit.
\end{proposition}
\begin{proof}
Since $\Delta/D$ is integral,
$\theta=\Delta/(c\sqrt q)=u\Delta/D$ is integral. The integral
basis of $E$ gives $\theta=(x+y\sqrt q)/2$ with the stated
parity. Total positivity of $\Delta$ gives $y>0$ and
$x^2<qy^2$. Its two dimension characters contribute trace $cqy$;
the other four codegrees exceed one. The signed Casimir trace is
an integer at most $T$, giving the displayed finite bound.

The global dimension is a $d$-number by
\cite[Theorem~1.2]{Ocode}. Here
\[
 \operatorname{Tr}(\Delta)=cqy,\quad
 \operatorname{N}(\Delta)=\frac{c^2qn}{4},\quad
 \operatorname{Tr}(1/\Delta)=\frac{4y}{cn}.
\]
The quadratic $d$-number divisibility gives $n\mid4qy^2$,
and the reciprocal-codegree identity gives the reciprocal equation.
The integral-basis coordinates of $\Delta$ are
$c(qy-x)/2$ and $cx$. Thus $\Delta/f_i$ is integral precisely
when $f_i$ divides their greatest common divisor. Divisibility
by original rational codegrees follows separately from induction.
Finally the signed and original matrices are congruent modulo two,
so their Casimir characteristic polynomials have the same reduction.
Every enumeration range is finite. The case $x=0$ is retained:
$K$ can be quadratic even when $\Delta$ itself is rational.
\end{proof}

\begin{theorem}\label{thm:r3-arithmetic-exclusions}
The rings $\RR{8}{6524}$, $\RR{8}{6525}$ and $\RR{8}{6999}$
have no complex fusion categorifications.
\end{theorem}
\begin{proof}
For the first two rings, $q=149$, $c=9$,
\[
 D=\frac{6705+549\sqrt{149}}2,\qquad T=6741.
\]
Their reductions at two have field degrees $(2,3,1)$ and
$(2,1,1,1,1)$, respectively. Proposition~\ref{prop:r3-local-isolation}(1)
applies. The finite spectrum test has $1\leq y\leq5$ and leaves
only
\[
 y=5,\qquad x=\pm61,\qquad (f_1,f_2,f_3,f_4)=(9,9,9,9).
\]
Both traces equal $6741$.

For $\RR{8}{6999}$, $q=41$, $c=18$,
\[
 D=3690+576\sqrt{41},\qquad T=7416.
\]
The local factors are $\mathbb F_8$, the dual numbers of rank two,
and $\mathbb F_2$. Proposition~\ref{prop:r3-local-isolation}(2)
applies. The bound is $1\leq y\leq10$; the only survivors are
\[
 y=10,\qquad x=\pm64,\qquad (f_1,f_2,f_3,f_4)=(9,9,9,9),
\]
again with total trace $T$.

In all three cases the signed Casimir trace is the sum of the
squares of the signed structure constants. Since each absolute
value is bounded by its original coefficient, equality of traces
forces equality of every absolute value. The remaining unrestricted
sign-only associativity systems are exhausted exactly. Every
solution is a dual-compatible sign gauge of the original tensor.
After this gauge, the lift fusion rules are the original positive
ones, and the recorded dimension-character tests leave only the
Galois orbit of the positive character. Any categorification would
therefore have a pseudo-unitary Galois conjugate.

The pseudo-unitary possibility is excluded by the exact
\texttt{global\_codegree\_moment} certificates, checked independently
on the full input tensors. For the first two rings the obstruction
is visible already in the integral span of simple dimensions.
They are $1,u-1,u,u,u,u$, where
$u=(37+3\sqrt{149})/2$ and $u^2-37u+7=0$.
The codegree $9$ would require a central simple of dimension
\[
 D/9=(61u-11)/3\notin\mathbb Z+\mathbb Z u,
\]
which cannot be the dimension of an integral restriction to the
original category. This also gives a direct check on the final
induction obstruction in those two cases.
\end{proof}

The tensor-only generator and exact verifier are
\filename{software/round3_arithmetic.py}; all local algebras,
surviving spectra, sign gauges, and complete general certificates
are recorded in \filename{research/round3/arithmetic/}.
The independent enumeration in
\filename{research/round3/models/arithmetic_audit.py} reproduces
the first pair of saturated spectra without an optimizer or the
generator's recursive divisor enumeration.

\subsection{Noncommutative signed spectra}
\label{sec:r3-noncommutative}

The preceding field-isolation argument has a different replacement
when the signed quotient is forced to be noncommutative. This follows
either from reduction modulo two or from a separate exhaustive
exclusion of the commutative signed sector. It is essential in this
case to keep the representation multiplicities.

\begin{proposition}\label{prop:r3-noncommutative-spectrum}
Let $R$ have rank six. Suppose either that its reduction modulo two
is noncommutative, or that every commutative signed quotient permitted
by sphericalization has been excluded by necessary conditions.
Suppose its two linear characters form a real quadratic orbit
over $E=\mathbb Q(\sqrt q)$, and let $D,D'$ be their formal
codegrees and $c$ that of its unique two-dimensional representation.
For any categorification there are $\Delta\in E$ and
$h\in\mathbb Z_{>2}$ satisfying
\[
 \frac1\Delta+\frac1{\Delta'}+\frac2h=1,
 \qquad \operatorname{Tr}_{E/\mathbb Q}(\Delta)+8h\leq T,
 \qquad T=\sum_{i,j,k}(N_{ij}^k)^2.
\]
Both embeddings of $\Delta$ are positive, $\Delta$ is a
$d$-number, and
\[
 \Delta/D,\quad \Delta/D',\quad \Delta/c,\quad \Delta/h
 \quad\hbox{belong to }\mathcal O_E.
\]
For fixed $R$ these conditions are a finite exact search.
\end{proposition}
\begin{proof}
The negative-sign quotient of sphericalization is now an oriented
signed algebra $A$. Its signed basis is orthonormal for the
coefficient-of-unit inner product, and duality gives adjoints of
left multiplication. Thus its regular representation is a finite
dimensional $*$-algebra with positive inner product and is
semisimple. In the first case parity preserves a nonzero commutator;
in the second case the commutative alternative is excluded separately.
In either case dimension six and the existence of a linear character give
\[
 A\otimes\mathbb C\cong\mathbb C\oplus\mathbb C\oplus M_2(\mathbb C).
\]
There are exactly two linear characters. The spherical dimension
character is real and dual invariant, and its field $K$ therefore
has degree at most two. Induction applied to the original characters
still gives $\Delta/f\in\mathcal O_K$. In particular $D\in K$,
so $K=E$. The signed linear codegrees are $\Delta,\Delta'$.

The unique two-dimensional representation is preserved by Galois
action. Its formal codegree $h$ is consequently rational and an
algebraic integer, hence an integer. The formal-codegree reciprocal
identity weights it by its representation dimension two. On the
four-dimensional regular matrix block the Casimir has eigenvalue
$2h$, so that block contributes $8h$ to the regular trace.
This proves both displayed identities, with the trace inequality
following from the coefficient bounds. Positivity of the two linear
codegrees and the reciprocal equation imply $h>2$.
Induction from the negative-sign character gives $\Delta/h$;
induction from the original characters gives the other quotients.
The usual global-dimension $d$-number theorem applies to $\Delta$.

For example, when $q\equiv1\pmod4$, write
$\Delta=(a+b\sqrt q)/2$, with $a\equiv b\pmod2$.
Then $0<a\leq T-24$ and $qb^2<a^2$. These bounds are finite.
For each pair $(a,b)$ the reciprocal equation determines
\[
 h=\frac{2}{1-a/\operatorname{N}(\Delta)}.
\]
Testing integrality, positivity, the four quotients, and the
quadratic $d$-number condition completes the exact necessary search.
\end{proof}

\begin{theorem}\label{thm:r3-noncommutative-exclusion}
The ring $\RR{6}{170}$ has no complex fusion categorification.
\end{theorem}
\begin{proof}
Its original codegrees are $148\pm24\sqrt{37}$ and $4$, with
representation dimensions $1,1,2$, and $T=328$.
The preceding finite system leaves exactly
\[
 \Delta=148\pm24\sqrt{37},\qquad h=4,
 \qquad \operatorname{Tr}(\Delta)+8h=328.
\]
Thus every signed coefficient has its original absolute value.
The complete oriented sign-only associativity enumeration has
$1631$ search nodes and $16$ solutions; all sixteen are
dual-compatible sign gauges of $R$. No commutativity or indicator
normalization is imposed. Under each gauge the two possible
dimension characters are Galois conjugates of the positive one,
so every categorification would have a pseudo-unitary Galois
conjugate.

Multiplicity-weighted central induction has $146$ candidate
restriction columns for the nontrivial linear character and $25$
for the two-dimensional representation. Their complete compatibility
search yields nineteen low frames. Every frame violates an ordinary
exact Gram or moment condition. The pseudo-unitary possibility is
therefore impossible, completing the general exclusion.
\end{proof}

The implementation is \filename{software/round3_noncommutative.py},
with arithmetic spectra, all oriented sign tables, the nineteen
induction branches, and replays in
\filename{research/round3/noncommutative/}.
For rings with commutative reduction modulo two, the first hypothesis
of Proposition~\ref{prop:r3-noncommutative-spectrum} is unavailable.
Their entire commutative signed sector must be treated before the
noncommutative spectrum calculation supplies a general reduction.

\begin{theorem}\label{thm:r3-noncommutative-sector-exclusions}
The rings $\RR{7}{200}$ and $\RR{8}{491}$ have no complex fusion
categorifications.
\end{theorem}
\begin{proof}
Both reductions modulo two are commutative. First impose commutativity
on the unrestricted oriented signed coefficient problem, identifying
opposite coefficient variables and intersecting their original parity
intervals. If an opposite original coefficient is zero, its identified
signed coefficient is forced to zero. This exhausts all commutative
signed quotients without an indicator normalization.

The complete enumerations have respectively $176$ and $208$ signed
tables, forming $11$ and $13$ dual-compatible sign-gauge orbits.
Their character algebras have respectively $27$ and $53$ components.
Each component has a zero dimension, unequal dimensions on dual
objects, or an exact norm-divisibility obstruction to induction.
The latter uses the original formal codegrees, whose quadratic
polynomial is $z^2-1560z+1872$, and the matrix-block codegree $12$.
The regular Casimir eigenvalue of that block is $24$; it must not
replace the formal codegree in the induction test. The recorded
norm obstructions occur at primes $3$ or $13$. Consequently no
categorical signed quotient is commutative.

Proposition~\ref{prop:r3-noncommutative-spectrum} now applies.
For both rings its finite system has exactly
\[
 \Delta=780\pm216\sqrt{13},\qquad h=12,\qquad
 \operatorname{Tr}(\Delta)+8h=1656=T.
\]
The resulting sign-only searches have $4927$ and $6463$ nodes,
respectively, and sixteen solutions each. Every solution is a
dual-compatible gauge of the original tensor. Thus any
categorification has a pseudo-unitary Galois conjugate. In each
ring the multiplicity-weighted induction problem has ten complete
low frames, all excluded by exact ordinary Gram or moment conditions.
This excludes the pseudo-unitary possibility and proves the claim.
\end{proof}

The separate commutative-sector verifier is
\filename{software/round3_commutative_sector.py}; its complete
certificate is included in each general certificate in
\filename{research/round3/noncommutative/}.
Both general replays also reconstruct the noncommutative arithmetic
and induction calculations. Removing the sector certificate is
explicitly rejected by the verifier.

\paragraph{Scope and controls.}
Each detector reconstructs its hypotheses from the full tensor.
The new family and arithmetic criteria are reapplied to the complete
census through multiplicity eight; exclusions of already resolved
rings are counted as confirmations. Exact replay disables numerical
optimizers and includes altered-certificate rejection and changes
of basis. All $54$ positively realized rings are controls, with
ineligibility distinguished from survival of an applicable system.
The C4 family includes the positive parameter-zero member in its
applicability scan, although the present irrational-dimension trace
engine does not apply to it. Its positive parameter-two member
survives the necessary system. These controls check scope; they do
not turn a surviving necessary condition into a categorical model.

\subsection{Dimension-block Fourier bounds and saturated twist support}
\label{subsec:dimension-block-support}
Let $K_1=\QQ(d_i)$, $K_0=\QQ(D)$ and $e=[K_1:K_0]$ for a spherical
dimension function. For a character $\chi$ of
$\operatorname{Gal}(K_0/\QQ)$, let $f_\chi$ be its primitive Dirichlet
conductor. Schopieray's dimension-block theorem
\cite[Theorem~5.1]{SchDimensionBlock} gives the additional necessary bound
\begin{equation}\label{eq:dimension-block-fourier-bound}
 \left|\sum_{\sigma\in\operatorname{Gal}(K_0/\QQ)}
       \frac{\chi(\sigma)}{\sigma(D)}\right|^2
 \le \frac1{e^2 f_\chi}.
\end{equation}
Its equality case supplies a restriction on individual twists, rather
than only on their averages. Here $S$ denotes the center's modular
$S$-matrix normalized by $D$, and $T$ is its matrix of actual twists.

\begin{proposition}[Saturated conductor support]
\label{prop:saturated-conductor-support}
Suppose equality holds in \eqref{eq:dimension-block-fourier-bound}.
Inflate $\chi$ to $G_1=\operatorname{Gal}(K_1/\QQ)$, and let
$w_\chi=\sum_{\sigma\in G_1}\chi(\sigma)e_{A_\sigma}$, where
$A_\sigma$ is the low central summand indexed by the dimension
character $\sigma(d)$. The cyclic subrepresentation generated by
$w_\chi$ in the canonical modular representation of $Z(\mathcal C)$
factors through $\mathrm{SL}_2(\mathbb Z/f_\chi\mathbb Z)$.
Consequently, for a central simple with restriction column $v$,
\begin{equation}\label{eq:saturated-column-support}
 L_\chi(v):=\sum_{\sigma\in G_1}\chi(\sigma)
       \sigma\!\left(\frac{\sum_i v_i d_i}{D}\right)\ne0
 \quad\Longrightarrow\quad \theta_Z^{f_\chi}=1.
\end{equation}
If several saturated characters give nonzero coefficients, the twist
order divides the greatest common divisor of their conductors.
\end{proposition}
\begin{proof}
We extract the equality information from the local compression
calculation in \cite[Lemma~4.1]{SchDimensionBlock}. Let
$\Gamma=\mathrm{SL}_2(\mathbb Z/\ell^a\mathbb Z)$, let $U$ be its
upper unipotent subgroup, and let $P$ project onto $U$-fixed vectors.
For a nontrivial torus character $\alpha$ of conductor $\ell^c$, the
universal permutation representation on primitive columns realizes
$\operatorname{Ind}_U^\Gamma\mathbf1$. Its $\alpha$-space of
$U$-invariants splits into the one-dimensional big-cell space and
the small-cell space. The compression $PsP$ maps the big-cell line
to the small-cell space and maps the latter back to a line. Both
nonzero singular values are $\ell^{-c/2}$. Hence every vector
attaining this norm belongs to the span of the big-cell vector and
the adjoint image of the opposite-character big-cell vector.

Writing a primitive column as $(p,q)$, these two vectors are, up to
nonzero scalar factors,
\[
 \alpha(q)[q\text{ is a unit}],\qquad
 \alpha(p)[q\equiv0\pmod{\ell^c}].
\]
Both are pullbacks of functions on primitive columns modulo $\ell^c$.
Their entire $\Gamma$-cyclic spans therefore factor through reduction
modulo $\ell^c$. Every irreducible representation with $U$-invariants
embeds equivariantly and isometrically, up to a scalar, into this
permutation representation. Thus an irreducible constituent containing
a nonzero norm-attaining vector itself factors through that reduction.

For the trivial local character the bound is $1$. Equality implies
that $sw$ is $U$-fixed as well as $w$. The vector $w$ is therefore fixed
by both upper and lower unipotents. These generate $\Gamma$, so its
cyclic representation is trivial. Decompose a general unitary
representation orthogonally into irreducibles, and then use the Chinese
remainder tensor decomposition at the prime divisors of its level.
Equality in the product norm bound forces every occupied local factor
to attain its own bound. The preceding conclusions show that a
norm-attaining vector generates a representation of level dividing
$f_\chi$.

The proof of the dimension-block bound estimates the matrix coefficient
$\langle w_{\bar\chi},PSPw_\chi\rangle$ by
$f_\chi^{-1/2}\|w_\chi\|^2$. Equality in the bound forces equality
in this compression-norm estimate and in Cauchy--Schwarz. The preceding
argument applies to $w_\chi$. In particular
$(T^{f_\chi}-I)Sw_\chi=0$. The known dimension-character rows give
$S_{Z,A_\sigma}=\sigma(d_Z/D)$, so the $Z$-coordinate of $Sw_\chi$
is precisely $L_\chi(v)$. This proves
\eqref{eq:saturated-column-support}. The greatest-common-divisor
assertion follows by intersecting root-of-unity groups.
\end{proof}

The exact audit \filename{research/round5/unification/dimension_block_bound.py}
enumerates every dimension-field character for the six unresolved rings.
All satisfy the bound. Their nontrivial equality conductors are
$20$ for $R_{6,4,22}$, $16$ for $R_{6,4,91}$, $15$ for $R_{6,4,629}$,
$12$ for $R_{6,5,1204}$, and $12,24$ for $R_{6,8,1660}$;
$R_{6,4,155}$ has strict inequalities. The receipt records all Fourier
values, including trivial characters, in exact cyclotomic coordinates.
The reusable implementation \filename{software/dimension_block_support.py}
computes \eqref{eq:saturated-column-support} from the dimensions and each
restriction column. When the forced roots lie in the retained field,
their complete power tuples filter the projected twist types exactly.

\begin{proposition}[Universal moments of a saturated real character]
\label{prop:saturated-modular-moments}
In Proposition~\ref{prop:saturated-conductor-support}, suppose $\chi$
is real and put $f=f_\chi$,
$\lambda=\sum_{\sigma\in G_1}\chi(\sigma)/\sigma(D)$.
Let $M=\operatorname{Ind}_{B}^{\mathrm{SL}_2(\mathbb Z/f\mathbb Z)}\chi$,
where $B$ is the upper triangular subgroup and its character is trivial
on unipotents and equals $\chi$ on the diagonal torus. Write $S_f,T_f$
for the standard unitary operators on $M$, and $P_\chi$ for projection
onto its unipotent-fixed, torus-$\chi$ space. Let the columns of $K$
span the subspace on which $P_\chi S_fP_\chi$ has eigenvalue $\lambda$.
For every exact matrix identity
\[
 \sum_{n=0}^{f-1}c_n K^*S_f^*T_f^nS_fK=0,
\]
the following necessary moment equation holds:
\begin{equation}\label{eq:saturated-modular-moments}
 \sum_{n=0}^{f-1}c_n\sum_Z|L_\chi(v_Z)|^2\theta_Z^n=0.
\end{equation}
\end{proposition}
\begin{proof}
By the preceding proposition, the cyclic representation generated by
$w_\chi$ factors through this finite group. It embeds equivariantly and
isometrically into $M$. Indeed, in an isotypic summand
$\pi\otimes\mathbb C^m$, the generating vector lies in
$\pi^{B,\chi}\otimes\mathbb C^m$; cyclicity forces
$m\le\dim\pi^{B,\chi}$. Frobenius reciprocity gives exactly this latter
multiplicity in $M$. Complete reducibility and invariant inner products
make the embedding isometric.

Since $\chi$ is real, equality in Cauchy--Schwarz in the preceding
proof says $P_\chi Sw_\chi=\lambda w_\chi$. Thus its embedded vector
lies in the column space of $K$. Applying the displayed matrix identity
to that vector gives
$\sum_n c_n\langle Sw_\chi,T^nSw_\chi\rangle=0$.
The simple-coordinate expansion of this inner product is exactly
\eqref{eq:saturated-modular-moments}.
\end{proof}
The signed permutation model on the projective line, its projectors,
equality eigenspace and complete matrix-relation kernel are constructed
by \filename{research/round5/unification/equality_modrep.py}.
This uses the full equality eigenspace, so it does not select an
unproved irreducible constituent or a particular vector within it.
There is also a direct spectral support test. If $E_r$ is the orthogonal
projector onto the $\zeta_f^r$-eigenspace of $T_f$, compute
\[
 K^*S_f^*E_rS_fK=
 \frac1f\sum_{n=0}^{f-1}\zeta_f^{-rn}K^*S_f^*T_f^nS_fK.
\]
If this matrix is zero, positivity implies $E_rS_fK=0$.
Thus $L_\chi(v_Z)\ne0$ forbids $\theta_Z=\zeta_f^r$.
This exact discrete Fourier calculation can remove phase variables
before the moment equations are solved.

\begin{corollary}[Different saturated conductors are orthogonal]
\label{cor:saturated-cross-moments}
Let $\chi,\psi$ be saturated characters with distinct primitive
conductors $f_\chi,f_\psi$. Then for every integer $n$,
\begin{equation}\label{eq:saturated-cross-moments}
 \sum_Z\overline{L_\chi(v_Z)}L_\psi(v_Z)\theta_Z^n=0.
\end{equation}
Only the exponents $0\le n<\gcd(f_\chi,f_\psi)$ are needed.
\end{corollary}
\begin{proof}
Every irreducible constituent in the cyclic span of $w_\chi$ contains
a nonzero unipotent-fixed vector of torus character $\chi$. Its local
congruence level is at most the conductor exponent, by
Proposition~\ref{prop:saturated-conductor-support}, and cannot be smaller:
otherwise this torus character would factor through a smaller modulus.
The constituent therefore has exact conductor level $f_\chi$.
Constituents belonging to distinct conductors cannot be isomorphic.
Orthogonal complete reducibility of the unitary modular representation
now makes the two cyclic subspaces orthogonal. In particular
$\langle Sw_\chi,T^nSw_\psi\rangle=0$, which is the displayed formula.
Every nonzero summand has both Fourier coefficients nonzero, so its
twist order divides both conductors. The expression is therefore
periodic with period their greatest common divisor.
\end{proof}

\subsection{Different ideals and a split quadratic dimension field}
\label{sec:round3-different}
The remaining dense quadratic profile requires a refinement that also controls
possible larger signed dimension fields.

\begin{theorem}[Different-ideal obstruction]\label{thm:signed-different}
Let $K$ be the dimension-character field of the commutative signed algebra
obtained by sphericalization of a commutative self-dual fusion ring.
Its global dimension $\Delta$ belongs to the different
$\mathfrak D_{K/\mathbb Q}$.
\end{theorem}
\begin{proof}
Write $d_i$ for the signed basis dimensions, and let $e_K\in K\otimes K$
be the separability idempotent. Orthogonality of distinct Galois-conjugate
characters gives
\[
 (\Delta\otimes1)e_K=\sum_i d_i\otimes d_i
                 \in\mathcal O_K\otimes\mathcal O_K.
\]
For $a$ in the codifferent, contract the second factor by
$x\mapsto\operatorname{Tr}(ax)$. The right side is integral, and the left
side is $\Delta a$. Thus $\Delta\mathfrak D^{-1}\subseteq\mathcal O_K$,
which is the assertion.
\end{proof}

For a local extension with ramification index $e$ and different exponent $d$,
\[
 \operatorname{Tr}(\mathfrak p^k)
       =p^{\lfloor(k+d)/e\rfloor}\mathbb Z_p.
\]
In particular, every tamely ramified odd prime divides the trace of the
different. A cyclic cubic extension of conductor $9$ has $e=3,d=4$ and
contributes $9$. We use these standard local trace formulas and the
conductor--discriminant formula; see Serre, \emph{Local Fields},
Chapters III and VI, and Milne, \emph{Class Field Theory}, Chapter V, Theorem~3.27.\footnote{The latter reference is available at
\url{https://www.jmilne.org/math/CourseNotes/CFT.pdf}.}

\begin{proposition}\label{prop:split-41-profile}
A commutative self-dual rank-six fusion ring with formal codegrees
\[
 9,9,9,9,\qquad 3690\pm576\sqrt{41},
\]
Casimir trace $7416$, and reduction
$\mathbb F_2^4\times\mathbb F_2[\epsilon]/(\epsilon^2)$ satisfies the following
necessary condition for categorifiability: its signed Casimir trace is $7416$.
Consequently every signed coefficient has its original absolute value.
\end{proposition}
\begin{proof}
Put $E=\mathbb Q(\sqrt{41})$. The dimension-field containment lemma gives
$E\subseteq K$, $\Delta/9\in\mathcal O_K$, and
\[
 \frac{\Delta}{18\sqrt{41}}\in\mathcal O_K,
 \qquad 738\mid\operatorname{Tr}_{K/\mathbb Q}(\Delta),
 \qquad \operatorname{Tr}_{K/\mathbb Q}(\Delta)\le7416.
\]
Indeed $3690+576\sqrt{41}=18\sqrt{41}(32+5\sqrt{41})$, and the last factor
is an integral unit of norm $-1$. Cyclotomicity makes $K/\mathbb Q$
abelian Galois, of degree $2$, $4$, or $6$.

If $[K:\mathbb Q]$ is $4$ or $6$, then $2$ splits completely in $K$.
The four etale rank-one blocks of the completed signed order have generic
fibre $\mathbb Q_2$; its only other block has rank two. A Galois number
field using a completion of degree two would require at least two such
local components, whereas only one rank-two block is available. The same
argument excludes ramification at $2$. It allows rational idempotents
inside the non-etale block and does not assume that block is maximal.

Theorem~\ref{thm:signed-different} and the trace bound restrict ramified
primes outside $\{3,41\}$ to $5$ and $7$, and forbid their simultaneous
occurrence, which would require trace at least $35\cdot738$.
The complete list of possible larger fields follows from elementary
Dirichlet-character arithmetic:
\begin{enumerate}
\item A biquadratic field $E\mathbb Q(\sqrt m)$ requires $m>1$ squarefree,
positive, prime to $41$, supported on $\{3,5,7\}$, and $m\equiv1\pmod8$.
The only such product is $105$, already forbidden by its primes $5$ and $7$.
\item A cyclic quartic character with square the quadratic character of
conductor $41$ is a quartic character modulo $41$ times a quadratic
character of positive odd discriminant $m$. Since $2^{10}\equiv-1\pmod{41}$,
splitting at $2$ requires $m\equiv5\pmod8$. The permitted products leave
exactly $m=5,21$. The conductors are $205,861$, and the possible dimension
traces are respectively $3690,7380$ and $5166$.
\item A degree-six field is $E L$ with $L$ cyclic cubic. Its conductor is
$7$, $9$, or $63$: neither $5$ nor $41$ admits a cubic character, and a
nontrivial cubic character at $3$ has conductor $9$. The prime $2$ is inert
in the conductor-$7$ and conductor-$9$ cubic fields. Only conductor $63$
remains, with dimension trace $5166$.
\end{enumerate}
There are no additional prime-power conductors: characters of order dividing
$4$ or $6$ are trivial on the higher prime-power unit groups at primes other
than $2$ or $3$; the relevant quadratic odd characters have squarefree conductor.

In the sextic case, coprimality of the discriminants gives
$\mathcal O_{EL}=\mathcal O_E\otimes\mathcal O_L$ and a product of differents.
Write
\[
 \Delta=18\sqrt{41}\,\frac{a+b\sqrt{41}}2,
 \qquad a,b\in\mathcal O_L,\quad a\equiv b\pmod2.
\]
Positivity at both embeddings of $E$ over each real embedding of $L$ makes
$b$ totally positive, with $\operatorname{Tr}_{L/\mathbb Q}(b)=7$.
At $7$, the cubic different is $\mathfrak p_7^2$, and $18$ is a unit.
The different theorem therefore forces $b\in\mathfrak p_7^2$. Hence
\[
 49\le\operatorname{Norm}(b)\le(7/3)^3<49,
\]
a contradiction.

The cyclic quartic cases require a short exact lattice enumeration. For
$M=5,21$, set
\[
 \alpha^2=\frac{M(41+5\sqrt{41})}{2},\qquad
 \beta^2=\frac{M(41-5\sqrt{41})}{2},\qquad
 \alpha\beta=2M\sqrt{41}.
\]
The defining polynomial is $X^4-41MX^2+164M^2$, and an integral basis is
\[
 1,\quad\frac{1+\sqrt{41}}2,\quad
 \frac{1+\sqrt{41}+2\alpha}4,\quad
 \frac{1-\sqrt{41}+2\beta}4.
\]
The checker verifies its integral multiplication, cyclic Galois action,
complete splitting at $2$, and discriminant $41^3M^2$. The preceding field
list and the conductor--discriminant formula prove maximality.
For $\theta=\Delta/(18\sqrt{41})=\sum_jz_jb_j$ and trace-dual basis $b_j^*$,
\[
 z_j=\operatorname{Tr}\left(\Delta\frac{b_j^*}{18\sqrt{41}}\right).
\]
At fixed positive trace $T$, this lies between $T$ times the smallest and
largest conjugates of $b_j^*/(18\sqrt{41})$. Exact algebraic sign comparisons
certify a finite integer box for every $z_j$. Imposing the trace and total
positivity leaves $28,216,16$ candidates at traces $3690,7380,5166$,
respectively. Every candidate fails a recorded $d$-number divisibility.
The supplementary certificate contains all $260$ candidates and the exact
coefficient witnessing each failure.

It remains to consider $K=E$. Write
$\Delta=18(41y+x\sqrt{41})/2$. Then $x\equiv y\pmod2$, $1\le y\le10$,
$x^2<41y^2$, and
\[
 41y^2-x^2\mid164y^2.
\]
The last condition is the quadratic $d$-number test. The four complementary
codegrees form either four rational integers, two rational integers and a
quadratic pair, or two quadratic pairs in $E$. Rational codegrees divide
the integral content of $\Delta$. A quadratic pair has polynomial
$X^2-SX+P$, with $S^2-4P=41r^2$, $P\mid\operatorname{Norm}(\Delta)$,
$P\mid S^2$, both roots greater than one, and both dividing $\Delta$
integrally. The total trace bounds $S$ and $r$. The reciprocal identity
$\sum_f1/f=1$ completes the finite test. Exactly seven dimension candidates
reach this stage; none admits a quadratic complementary pair. The only
complete spectra have $y=10$, $x=\pm64$, and four codegrees $9$.
Their total trace is $7416$, proving the proposition.
\end{proof}

The full tensor of \texttt{R6\_8\_6526} satisfies these hypotheses.
Exact saturated signed associativity leaves one table up to the certified
sign gauges, equivalent to the original tensor. Every non-FP original
character orbit fails the arithmetic dimension-character tests. Thus a
categorification would have a pseudo-unitary Galois conjugate, which the
existing global-codegree moment certificate excludes. The resulting
\emph{general} exclusion is implemented in
\filename{software/round3_split_arithmetic.py}, with canonical certificate
\filename{research/round3/arithmetic/R6_8_6526_general.json}.

\subsection{The dimension order and signed central traces}\label{sec:dimension-order}
The positive-dimension induction system has two extensions that are useful
when an exhaustive sphericalization search leaves a non-positive
spherical character. The first sharpens an arithmetic condition without
adding any central variables. The second retains the algebraic part of
the same trace system, but discards every inequality that depends on
positivity of categorical dimensions.

\subsubsection{The actual order of simple dimensions}
Let $d=(1,d_1,\ldots,d_{r-1})$ be a dimension character under
consideration, and let $K=\QQ(d_i)$ and $\delta=\sum_i d_i d_{i^*}$.
The finitely generated group
\[
 \mathcal O_d=\sum_i\ZZ d_i\subset K
\]
is an order: it contains $1$, consists of algebraic integers, and is
closed under multiplication because the $d_i$ satisfy the multiplication
equations with integer coefficients. It need not be the maximal order
of $K$.

\begin{proposition}[Dimension-order obstruction]\label{prop:dimension-order}
If a fusion ring $R$ has a spherical categorification with simple
dimensions $d_i$, then for every formal codegree $f$ of $R$,
\begin{equation}\label{eq:dimension-order}
 f\in K,\qquad \delta/f\in\mathcal O_d.
\end{equation}
The same conclusion holds for the dimensions of selected lifts in the
canonical sphericalization of an arbitrary categorification of $R$.
Thus the order-membership test may be applied to each dimension-character
branch of a complete signed-sphericalization search.
\end{proposition}
\begin{proof}
In the spherical case the central simple $A_E$ attached to an irreducible
representation of $R$ with codegree $f$ has dimension $\delta/f$, by
\cite[Theorem~2.13]{Ostrik3}. Its forgetful image is an object of the
original category. Its dimension is therefore a nonnegative integral
linear combination of the $d_i$, in particular an element of
$\mathcal O_d$. Since $\delta$ is nonzero, this also puts $f$ in $K$.

For an arbitrary categorification use its sphericalization, of dimension
$2\delta$. An irreducible representation of $R$ lifts through the
forgetful quotient to one with formal codegree $2f$; both selected and
opposite basis lifts contribute to that codegree. The same central
induction theorem gives dimension $(2\delta)/(2f)=\delta/f$.
All simple dimensions of the sphericalization are among $\{d_i,-d_i\}$,
so forgetting this central object gives an integral linear combination
of the selected dimensions. This proves the assertion without assuming
that the selected lifts themselves form a fusion subcategory.
\end{proof}

The test is finite and exact. Express the $d_i$ in a rational power basis
of $K$, clear one common denominator, and put their coordinate lattice
in Hermite normal form. Membership of $\delta/f$ is then an integer
linear-algebra question. A negative certificate can equivalently give a
rational linear functional $\ell:K\to\QQ$ satisfying
\[
 \ell(d_i)\in\ZZ\quad\hbox{for all }i,
 \qquad \ell(\delta/f)\notin\ZZ.
\]
The checker verifies the defining polynomial of $f$, reconstructs the
lattice and the functional, and checks both displayed conditions. A
failure on one character component is not a noncategorification proof
unless all possible sphericalization and character branches are closed.
In particular a failed character of a positively realized ring cannot
be promoted to an exclusion of that ring.

For illustration, the complete certificate for $\RR{4}{375}$ has four
oriented signed tables, all gauges of the original tensor. Its nonzero
spherical components have order $\ZZ[\sqrt5]$ and dimensions
\[
 \delta=70+30\sqrt5\quad\hbox{or}\quad\delta=30-10\sqrt5.
\]
The codegree $4$ gives quotients outside that order in both cases.
The supplied functional certificates, not this identifier or this
displayed calculation, are the inputs to the automated verification.

\subsubsection{The algebraic trace system without positive dimensions}
Suppose a spherical character takes values in a real field $K$ of degree
at most two, has nonzero simple dimensions, and satisfies
$d_i=d_{i^*}$. Write $\delta=\sum_i d_i^2$. The restriction vectors of
central simples remain nonnegative integral vectors, even when some
categorical dimensions are negative. Consequently the Gram identity
\eqref{eq:fullgram}, its finite entry bounds, and the formal-codegree
low-column equalities remain valid. They do not use an inequality
between categorical dimensions.

For a commutative input let $\chi_a$ be a character label of a low simple
$Z_a$ and put $\beta_a=\delta/f_a$. If $v_W$ is a candidate restriction
vector, define
\[
 d_W=d\cdot v_W,\qquad
 H_{aW}=\beta_a\chi_a(v_W),\qquad
 \Lambda_W=(H_{aW}/d_W)_a.
\]
Assume that the full ring character table is defined over $K$. The
character correspondence of Theorem~\ref{thm:character-induction} gives
$H_{aW}=S_{Z_a,W}$, with $S$ unnormalized. The following formulation
uses no absolute value of $d_W$.

\begin{theorem}[Signed central-induction obstruction]\label{thm:signed-traces}
A spherical categorification with the specified character supplies a
nonnegative integral solution of the following finite necessary system.
The unknowns count high central restriction vectors and their joint
relative twist-trace types. They satisfy the Gram equations, the signed
identities
\begin{equation}\label{eq:signed-indicators}
 \sum_W(v_W)_i d_W\theta_W^j=\delta\nu_j(b_i),
\end{equation}
and the cyclic-operator bounds on the right-hand indicators. Low
columns have twist one and dimension $\beta_a$. They admit at least
one complete character labelling satisfying modular reciprocity.

All coordinates of $\Lambda_W$ are algebraic integers. Their joint
Galois types have balanced multiplicities, separately for each relative
twist-trace type. An absent Galois type forbids the corresponding
restriction vector. Finally, if $P$ is the normalized trace from a
common cyclotomic field to $K$, then
\begin{equation}\label{eq:signed-low-modular}
 \sum_W P(\theta_W)H_{aW}H_{bW}=\delta S_{Z_a,Z_b}
\end{equation}
for every pair of low labels.

Exhaustive failure of the system for all bounded low columns, all
compatible character labellings, and all permitted second indicators
excludes the specified spherical character. Exact rational Farkas
certificates suffice for failure of the nonnegative rational relaxation.
\end{theorem}
\begin{proof}
Adjunction proves the Gram identity with nonnegative restriction
multiplicities. Central induction supplies the low dimensions and their
unit multiplicities exactly as in \eqref{eq:weightedlow}; neither
statement assumes pseudo-unitarity. The trace formula
\cite{NS} proves \eqref{eq:signed-indicators} with the actual, possibly
negative, dimensions. Bounds obtained as traces of finite-order cyclic
operators on ordinary complex vector spaces do not assume positivity
of categorical dimensions. Only these algebraic bounds, and not the
positive low-mass inequalities, are retained here.

The finite relative-trace list is justified by the exponent-prime
support theorem \cite{BNRW} and the same relative-trace calculation
used in \eqref{eq:orderbox}, with $\delta$ in place of $D$. The high
columns are bounded by the positive integral Gram matrix, not by the
sign of $d_W$. Their additional algebraic-integrality and $d$-number
conditions follow from the spherical modular center as in the proof of
Theorem~\ref{thm:balanced}. Columns of dimension zero are impossible.

Each $\Lambda_W$ is the restriction of the simultaneous character of
integer fusion matrices of the center to the labelled low objects.
It therefore consists of algebraic integers, and its full Galois orbit
occurs with the corresponding multiplicities. For the normalized
center representation $S/\delta,T$, the permutation of columns sends
$\theta_W$ to $\sigma^2(\theta_W)$; see
\cite[Theorem~2.9]{BNRW}. A square fixes the real field $K$ of degree
at most two, so normalized relative traces of every twist power are
unchanged. This proves the asserted balances of joint eigenvalue types.
It does not identify $d_W$ with its absolute value, nor impose the
positive dimension-orbit rule \eqref{eq:augmented-galois}.

The center has Gauss sum $\delta$. Its modular relation, together with
the twist-one property of low objects, is
$\sum_W\theta_W S_{Z_a,W}S_{W,Z_b}=\delta S_{Z_a,Z_b}$.
All coefficients other than the twists lie in $K$. Taking $P$ proves
\eqref{eq:signed-low-modular}. The actual central simples give integer
counts satisfying all the displayed equations. Expanding them in a
rational basis of $K$ and relaxing to nonnegative rational counts can
only enlarge the solution set. Farkas' lemma therefore proves the
last assertion.
\end{proof}

\paragraph{Exact scope of the implementation.}
The routine first checks the complete signed character table, its
multiplication equations, duality, nonzero dimensions and field degree.
It enumerates low columns using Gram bounds and exact dimension
equalities, without calling a positive-dimension low-mass test. It then
uses signed versions of the inherited indicator rows and adds joint
spectrum and low modular rows. In particular it does not reuse the
absolute-value dimension balances of
Proposition~\ref{prop:modular-frontier}. Only the second-indicator rows depend on the chosen signs, and each such row depends on one sign. The checker reconstructs the all-minus and all-plus endpoint systems and selects the exact normalized row for each sign choice. Independent mixed-sign reconstructions check this optimization; every separating-vector inequality is still verified. The second indicators of
self-dual objects are both allowed unless an independently checked
condition forces their sign.

For $\RR{4}{446}$ the remaining non-positive component is
\[
 d=(1,1,-1,-1,2-\sqrt6,2+\sqrt6),\qquad\delta=24.
\]
There are $290$ low matrices. Reciprocity leaves two, each with two
character labellings; all $32$ choices of second indicators are retained
for each labelling. The $128$ exact separation certificates exclude
these systems. The positive character component has the preceding
character-induction obstruction. The full oriented search has sixteen
gauge tables, so the assembled certificate excludes general
categorification, not just this displayed spherical structure.

\subsubsection{Application and structure-specific conclusions}
Component-level rejection and general nonexistence are different
conclusions. For example, rejecting one character of $\RR{4}{174}$ does
not exclude that ring. Complete proof trees for the order and signed
framework include
\[
 \RR{4}{375},\quad\RR{4}{446},\quad
 \RR{5}{1507},\quad\RR{5}{1508},\quad\RR{5}{1582},\quad
 \RR{6}{2767},\quad\RR{7}{1907}.
\]
For the last five, shorter norm and moment certificates close the
exhaustive signed alternatives, and these are retained rather than
credited to a stronger criterion unnecessarily. The finite applications
are specified by full tensors and exact proof trees.

The shorter pseudo-unitary certificates for $\RR{7}{3582}$ and
$\RR{7}{3812}$ are intermediate ingredients. Their unrestricted
existence branches are closed by the ramified signed reductions and
Casimir trace criteria above; both rings are generally excluded.

\subsection{Sphericality and Galois reduction at multiplicity three}
\label{sec:frontier-spherical-reduction}
\begin{proposition}\label{prop:frontier-galois-pu}
Every complex fusion category with based ring $\RR{3}{215}$ admits a
spherical structure and has a pseudo-unitary Galois conjugate.
Consequently general categorifiability of this ring is equivalent to
pseudo-unitary categorifiability. This reduction alone asserts neither
existence nor nonexistence; the final obstruction is proved in
Theorem~\ref{thm:frontier-noncategorifiable}.
\end{proposition}
\begin{proof}
Apply the exhaustive oriented signed-sphericalization construction of
Section~\ref{sec:signed}, using every coefficient in its full parity
interval and imposing no positive-indicator or modulo-four normalization.
There are $1480$ associative signed quotients. The $32$ unit-preserving
sign gauges partition them into $63$ orbits; one orbit consists of the
$32$ gauges of the original fusion tensor. On representatives of the
other $62$ orbits there are $271$ rational character components.
Of these, $228$ have a zero simple dimension and are impossible.
For each of the other $43$, the norm of the prospective global
dimension has $17$-adic valuation zero. The original ring has a formal
codegree with minimal polynomial $t^2-272t+1088$, whose norm has
$17$-adic valuation one. Thus the signed-sphericalization
norm-divisibility condition rejects every remaining component. No
nongauge quotient can occur for an actual category.

For completeness, a gauge quotient gives sphericality of the original
category, not merely compatibility with spherical dimensions. If
$T_{ij}^k=\eta_i\eta_j\eta_kN_{ij}^k$, replace the chosen lift of $X_i$
in the sphericalization by its tensor product with the sign object
when $\eta_i=-1$. The new signed quotient has $T=N$. The sum and
difference of the multiplicities of the two lifts of $X_k$ in a tensor
product are then both $N_{ij}^k$, so the opposite lift has multiplicity
zero. These chosen lifts therefore generate a full tensor subcategory
of the sphericalization which projects tensor-equivalently onto the
original category. Transport its spherical structure along that
equivalence.

The four rational characters of the original ring have a zero entry,
so none is a spherical dimension character. The only remaining
character orbit is
\[
 (1,1,d,d,d-1,d+1),\qquad d^2-8d-1=0.
\]
One of its two real embeddings is the Frobenius--Perron character,
$d=4+\sqrt{17}$. Galois conjugation to that embedding gives a spherical
category with positive Frobenius--Perron dimensions, hence a
pseudo-unitary category. Conversely a pseudo-unitary categorification
is a complex fusion categorification.

The complete unrestricted signed enumeration has $274655$ search nodes
and is frozen in
\filename{research/new_frontier/frontier_spherical_reduction.json}.
The companion Python script recomputes all $1480$ quotients and the
complete sign orbits, verifies the original gauge orbit, checks all
$271$ nongauge character components and their arithmetic obstructions,
and reconstructs the original dimension-character orbit using exact
arithmetic. The older normalized twelve-table enumeration is retained
separately as a regression control; it is not needed for this proof.

\end{proof}



